\section{Experimental Setup}
\label{sec:experimental-setup}

\paragraph{Models.}
We evaluate on two 7B-parameter models sharing the Qwen2 backbone (28 layers, 28 heads, 4-group GQA): \textbf{DeepSeek-R1-Distill-Qwen-7B}~\cite{deepseek-r1-2025}, distilled from DeepSeek-R1 with explicit chain-of-thought traces yielding pronounced phase structure; and \textbf{Qwen2.5-7B-Instruct}~\cite{qwen25-2024}, with standard instruction tuning, providing a cross-training-paradigm replication target. Both are loaded in \texttt{bfloat16} with eager attention.\footnote{An earlier \texttt{float16} configuration produced sporadic NaN values; \texttt{bfloat16} eliminated all occurrences.}

\paragraph{Dataset.}
We evaluate all 500 problems from MATH-500~\cite{math500-neurips2021}. Each problem is decoded greedily (\texttt{max\_new\_tokens}$=$2048). Problems with degenerate phase assignments (null low or high phase) are excluded, yielding effective sample sizes of $n{=}496$ for DeepSeek (4 excluded) and $n{=}475$ for Qwen (25 excluded). Qwen exhibits a higher mid-phase null rate (45.2\% vs.\ 10.2\% for DeepSeek) due to lower temporal autocorrelation in its attention sparsity patterns (lag-1 $r{=}0.718$ vs.\ $0.894$, $d{=}-1.68$); this does not affect the primary gradient metric $\Delta = \mathrm{II}_{\textit{low}} - \mathrm{II}_{\textit{high}}$, which requires only low and high phases. Sensitivity analysis across $\texttt{min\_phase\_length} \in \{5, 10, 20\}$ supports robustness (Appendix~\ref{app:mpl-sensitivity}).

\paragraph{Sparsity configuration.}
We apply symmetric compound sparsity at $s \in \{0.15, 0.20, 0.25, 0.30\}$: at each level, the head pruning rate $p{=}s$ (fraction of heads removed per layer) equals the KV cache eviction rate $e{=}s$ (fraction of positions evicted). We compute perplexity under four conditions: dense, attention-only ($p{=}s, e{=}0$), KV-only ($p{=}0, e{=}s$), and compound ($p{=}s, e{=}s$). We compare \emph{importance-based} eviction (position-normalized cumulative attention score, Eq.~\ref{eq:kv-importance}; BOS always retained) and \emph{random} eviction (uniform). Heads are pruned per layer by Frobenius norm; under random eviction, both heads and positions are selected uniformly.

\paragraph{Phase segmentation and metrics.}
Each sequence is segmented into three phases using the quantile-based method (\S\ref{par:phase-segmentation}), with $\texttt{min\_phase\_length}{=}10$. We compute segment-level perplexity and II (Eq.~\ref{eq:ii}), reporting medians due to heavy right-skew (mean/median ratios up to 43:1). The amplification ratio $G$ (Eq.~\ref{eq:gradient}) is reported in Appendix~\ref{app:g-metric}. Robustness to segmentation granularity (quartile, quintile; Appendix~\ref{app:seg-robust}) and minimum phase length ($5$, $10$, $20$ tokens; Appendix~\ref{app:mpl-sensitivity}) is verified.

\paragraph{Segmentation validation.}
Quantile-based segmentation shows strong agreement with model-intrinsic \texttt{<think>} boundaries on DeepSeek (Cohen's $\kappa{=}0.88$, token-level accuracy $94.2\%$, AUC-ROC${=}0.98$; $n{=}48$, GSM8K): $57.8\%$ of \texttt{<think>} tokens fall in the low-sparsity phase and $62.7\%$ of answer tokens fall in the high-sparsity phase, confirming that attention-sparsity phases recover genuine reasoning structure without requiring explicit delimiters. Boundary proximity is tight (median 1 token; $81.2\%$ within 10 tokens). Appendix~\ref{app:think-validation} provides full details.

\paragraph{Statistical analysis.}
\label{par:stats}
All reported confidence intervals use nonparametric bootstrap ($B{=}10{,}000$, seed$=$42)~\cite{efron-1993}. Significance tests use Wilcoxon signed-rank (paired within-model) or Mann-Whitney $U$ (cross-training-paradigm), with Bonferroni correction at family-wise $k{=}4$ (per model$\times$strategy combination, 4 sparsity levels testing one dose-response hypothesis): each family tests a single directional question (``does $\Delta$ increase with $s$?'') under fixed model and strategy, making $k{=}4$ the natural grouping. Benjamini-Hochberg FDR correction across all 44 tests is reported as a robustness check in Appendix~\ref{app:bh-fdr}; all Bonferroni-significant results remain significant under BH-FDR. Dose-response monotonicity is supported by the Jonckheere-Terpstra test ($z{=}6.2$--$33.7$, all $p < 10^{-9}$), which uses per-problem values (${\sim}496$ per group) rather than aggregated medians.\footnote{Spearman rank correlation on the four aggregate medians ($\rho{=}1.0$, $n{=}4$; interpret with caution given the small number of levels) is consistent; we rely on the JT test as primary evidence.} Scaling from $n{=}50$ (profiling) to $n{=}500$ narrowed bootstrap CI widths by $2.5$--$5.9{\times}$ across all 16 conditions, showing precision gains. The paired median difference $\Delta$ is the primary metric; $G$ (amplification ratio, Eq.~\ref{eq:gradient}) is reported in Appendix~\ref{app:g-metric} as an auxiliary illustrative ratio.\footnote{Experiments on 4$\times$ professional GPUs with 48\,GB memory.}
